\section{Système temporel dodécaphasique orienté du TPK}
\label{sec:ciclo-dodecafasico-tpk-rev7}
\index{TPK!système dodécaphasique}\index{périodicité d'ordre 108}

\HMTClaim{HMT-R-TPK-GAMMA108}

\subsection{Opérateur de sélection de route et transport vers l'observable dimensionnelle}

La fonction ternaire de phase opératoire est
\[
M_{\rm ph}(r)=
\begin{cases}
 +1,&r\in\{1,4,7\},\\
 -1,&r\in\{2,5,8\},\\
 0,&r\in\{3,6,9\}.
\end{cases}
\]
Ses trois valeurs sont les entrées scalaires avec lesquelles
\(\operatorname{Sel}_t\) active, respectivement, le curseur additif, le
curseur multiplicatif ou le compteur neutre. L'échantillonnage stroboscopique est
\[
 \operatorname{Obs}_3(a)=(a_3,a_6,a_9,\ldots).
\]
La sélection \(\operatorname{Sel}_t[\,M_{\rm ph}(g(t))\,]\) décide quelle
évaluation opère ; le transport cardinal déplace l'état ; et
\(\operatorname{Obs}_3\) sélectionne les instants qui sont publiés.

Le transport cardinal agit par quatre translations
\[
 T_d(i,j)=(i,j)+\nu_d\pmod9,
\]
\[
 \nu_N=(-1,0),\qquad \nu_E=(0,1),\qquad
 \nu_S=(1,0),\qquad \nu_O=(0,-1).
\]
Les curseurs additif et multiplicatif conservent des positions et des orientations
séparées. La coordonnée TRIT, la fonction \(M_{\rm ph}\), l'opérateur
\(\operatorname{Sel}_t\), les quatre translations et
\(\operatorname{Obs}_3\) sont donc des objets de types distincts.

\subsection{Calendrier et registre dodécaphasiques}

Le noyau topologique de phase ne s'arrête pas à la mise à jour d'une position
sur \(\mathbb Z_9^2\). Son état temporel minimal conserve l'ordre des
événements, l'orientation, le retour et la mémoire des transitions. Pour
une orbite de \(108\) pas, on introduit la partition ordonnée
\[
 E_m=\{9(m-1)+1,\ldots,9m\},
 \qquad m=1,\ldots,12,
\]
et l'on note son support ordonné
\[
 \Gamma_{108}=E_1\sqcup\cdots\sqcup E_{12},
 \qquad 108=12\cdot9.
\]
De manière équivalente,
\[
 \Theta_{108}=\mathbb Z/108\mathbb Z,\qquad
 \beta(\theta)=
 \left[\left\lfloor\frac{\widetilde\theta}{9}\right\rfloor\right]_{12},
 \qquad
 r(\theta)=1+(\theta\bmod9).
\]
La coordonnée \(\beta\) sélectionne l'un des douze secteurs et \(r\) localise la
phase dans sa fenêtre nonadique.

\begin{proposition}[Extension centrale non scindée de l'horloge observable]
\label{prop:extension-central-108-no-escindida-rev8}
En écrivant \(C_n=\mathbb Z/n\mathbb Z\), le calendrier observable forme la
suite exacte
\[
 0\longrightarrow C_{12}
 \xrightarrow{\ \iota\ }C_{108}
 \xrightarrow{\ p\ }C_9
 \longrightarrow0,
 \qquad
 \iota([b]_{12})=[9b]_{108},\qquad
 p([\theta]_{108})=[\theta]_9.
\]
Pour la section ensembliste \(s([r]_9)=[r]_{108}\), avec
\(0\leq r<9\), la couture est donnée par le cocycle de retenue
\[
 c(r,r')=
 \left[\left\lfloor\frac{r+r'}{9}\right\rfloor\right]_{12},
 \qquad
 s(r)+s(r')=s([r+r']_9)+\iota\bigl(c(r,r')\bigr).
\]
L'extension ne se scinde pas. Pour l'action triviale,
\[
 H^2(C_9,C_{12})
 \simeq C_{12}/9C_{12}
 \simeq C_3,
\]
et \(c\) représente la classe génératrice \([1]\).
\end{proposition}

\begin{proof}
L'image de \(\iota\) est le sous-groupe des multiples de neuf dans
\(C_{108}\), qui coïncide avec \(\ker p\) ; d'où l'exactitude. L'identité
du cocycle est la division euclidienne de \(r+r'\) par neuf. Si l'extension
centrale se scindait, le groupe médian serait isomorphe à
\(C_{12}\times C_9\). Cependant,
\[
 \exp(C_{108})=108,\qquad
 \exp(C_{12}\times C_9)=\operatorname{mcm}(12,9)=36,
\]
ce qui l'empêche. L'identification cohomologique pour un groupe cyclique à
action triviale donne \(H^2(C_9,C_{12})\simeq C_{12}/9C_{12}\), et la retenue
unitaire se projette sur le générateur de ce quotient.
\end{proof}

La non-scission type le retour : neuf itérations ferment la coordonnée dans
\(C_9\) et translatent d'une unité dans \(C_{12}\) ; cent huit itérations ferment les deux
coordonnées de cette horloge. Le calendrier \(C_{108}\) est la projection
temporelle de l'état enrichi \(\mathcal X_{\rm TPK}^{\rm enr}\) :
orientation, feuillet, frontière, cylindre, généalogie et mémoire profonde
restent dans ce dernier et peuvent
distinguer des histoires ayant le même calendrier observable.

Si \(g_0(\theta)=\theta\bmod9\), alors
\[
 g_0(\theta+9)=g_0(\theta),\qquad
 \beta(\theta+9)=\beta(\theta)+1\pmod{12}.
\]
Le témoin
\[
 (g_0,r,\beta,30^\circ\beta)(49)=(4,5,5,150^\circ),
\]
\[
 (g_0,r,\beta,30^\circ\beta)(58)=(4,5,6,180^\circ)
\]
montre exactement ce qui revient et ce qui conserve la mémoire : la phase nonadique et
la position interne coïncident, tandis que le secteur dodécaphasique avance d'une unité.
Le calibre régional \(\operatorname{diag}_c\) des régions de \(\pi\) n'est
pas cette coordonnée \(\beta\). Il n'existe pas non plus ici de morphisme
\(30^\circ\beta\to h\in\langle90,120\rangle\) ; en particulier,
\(150\notin\langle90,120\rangle\). La paire \(49/58\) est un témoin de retenue,
non un sélecteur régional ou spectral.

Pour une histoire enrichie \(x\) de TPK, soient \(\tau_m\) la sous-route
ordonnée sur \(E_m\), \(\sigma_m\) son orientation, \(\kappa_m\) la retenue
qui franchit la frontière du bloc et \(\Lambda_m\) l'état accumulé du
registre. L'objet
\[
 \mathcal R_{12}
 =\bigl((E_m,\tau_m,\sigma_m,\kappa_m,\Lambda_m)\bigr)_{m=1}^{12}
\]
est le système temporel dodécaphasique orienté du TPK. Il ne constitue pas une quatrième
théorie ajoutée à APP, TRIT et TPK : c'est la forme temporelle interne avec laquelle
TPK ordonne une orbite, conserve ses retours et publie un registre
réversible. On distingue les projections canoniques
\[
 \mathcal X_{\rm TPK}^{\rm enr}
 \longrightarrow \mathcal R_{12}
 \longrightarrow \Theta_{108}.
\]
La première conserve \((\tau_m,\sigma_m,\kappa_m,\Lambda_m)\) ; la seconde
publie uniquement la suite observable et ne peut pas reconstruire à elle seule
l'état signé. Le registre enrichi est noté \(\mathcal R_{12}\) et la
rotation des secteurs est notée \(C_{12}\).

La fermeture observable du calendrier satisfait
\[
 F_{\rm obs}^{108}=I.
\]
Les positions et les orientations sont déjà restaurées après \(54\) itérations,
tandis que \(\mathcal R_{12}\) distingue la position temporelle au sein des
douze secteurs. Ici \(F_{\rm obs}\) est la permutation observable, tandis que
\(\mathcal G_9\) est l'élagage au sein de la relation d'extension et de frontière.
La monodromie est la propriété conjointe
\[
 g(K+9)=g(K),\qquad v_{K+9}\ne v_K,
\]
non le nom isolé de l'élagage. La fermeture chronologique restitue le calendrier
observable ; le retour nonadique conserve la phase et modifie le préfixe, le cylindre
et la mémoire.

On distingue également deux mots de douze composantes :
\[
 U_{12}^{\rm cat}=U_6\Vert U_6,
 \qquad
 U_{12}^{\rm sgn}=\mathcal H_{12}(K).
\]
Le premier concatène deux émissions de longueur six. Le second est une
coordonnée signée de l'état enrichi et reste lié de façon réversible à
\(K\) par la transformation dodécaphasique \(\mathcal H_{12}\). Ni leurs domaines
ni leurs fonctions ne sont interchangeables ; en particulier, le mot concaténé
ne construit pas à lui seul la coordonnée signée.

\begin{definition}[Registre entier des douze événements]
\label{def:registro-dodecafasico-rev7}
Pour une route de chiffres \(d_1,\ldots,d_{108}\), soit
\[
 \mathcal D_{\rm evt}=\{2,4,6,8\},\qquad
 \Sigma_{12}=(+,+,+,-,-,-,+,+,+,-,-,-).
\]
Son mot orienté d'événements est
\[
 e_m=\Sigma_{12}(m)\,
 \#\{t\in E_{m+1}:d_t\in\mathcal D_{\rm evt}\},
 \qquad 0\le m<12.
\]
Les deux demi-couronnes sont appariées au moyen de
\[
 p_i=e_i+e_{i+6},\qquad
 a_i=e_i-e_{i+6},\qquad 0\le i<6.
\]
On définit le comptage brut, les cinq différences par rapport à la sixième paire et
les six antisymétries par
\[
 s_C=\sum_{i=0}^{5}p_i,\qquad
 M_i=p_i-p_5\quad(0\le i<5),\qquad
 A_6=(a_0,\ldots,a_5).
\]
Le registre entier est
\[
 \mathscr L_{12}(e)
 =\bigl(s_C,M_0,\ldots,M_4,a_0,\ldots,a_5\bigr).
\]
\end{definition}

\begin{theorem}[Coordonnées exactes de la mémoire dodécaphasique]
\label{thm:memoria-1-5-6}
L'application \(e\mapsto\mathscr L_{12}(e)\) est injective sur son image.
L'inverse est donné par
\[
 p_5=\frac{s_C-\sum_{i=0}^{4}M_i}{6},\qquad
 p_i=p_5+M_i\quad(0\le i<5),
\]
\[
 e_i=\frac{p_i+a_i}{2},\qquad
 e_{i+6}=\frac{p_i-a_i}{2}.
\]
En particulier, la distribution des coordonnées est
\[
 \boxed{1+5+6=12}.
\]
\end{theorem}

\begin{proof}
Par définition,
\[
 s_C=\sum_{i=0}^{4}(p_5+M_i)+p_5
 =6p_5+\sum_{i=0}^{4}M_i.
\]
Cela détermine \(p_5\) puis tous les \(p_i\). Additionner et soustraire les
équations \(p_i=e_i+e_{i+6}\) et \(a_i=e_i-e_{i+6}\) restitue les douze
événements. La divisibilité par six et les six conditions de parité
caractérisent l'image entière ; ni une masse ni une unité n'interviennent.
\end{proof}

La division par six ne s'applique pas pendant le comptage et n'est pas un coefficient
primitif. On forme d'abord le comptage orienté brut \(s_C\) ; ensuite, on
reconstruit le registre des douze événements ; ce n'est qu'alors qu'apparaît
\(s_C/6\) dans le caractère. De manière équivalente,
\[
 W_\pm=6n_A\pm s_C
\]
sont entiers et la factorisation réticulaire possède la forme normale de Smith
\((1,12)\). Cette antériorité sera obligatoire dans la loi des masses.

Les entiers \(90\) et \(120\) réapparaissent plus loin dans six constructions
de types distincts. Pour empêcher que le nom du cardinal ne se substitue à
l'objet, on fixe dès maintenant la convention suivante :

\begin{enumerate}
 \item \(\mathcal E_{\rm dod}=(D_3,D_4,Q):\mathbb Z^{12}\to\mathbb Z^{25}\),
       extracteur entier qui reconstruit \(K\) ;
 \item \(D_{\rm raw}^{90,120}(q)\), différence diagnostique de deux séries dans
       une carte fixée ;
 \item \(\iota_{6,8}:\mathcal F_6\hookrightarrow\mathcal F_8\), inclusion
       combinatoire de cardinalités \(90\) et \(120\) ;
 \item \(K_{6\to8}=12\Delta S_{90}-\Delta S_{120}\), fonctionnelle angulaire
       définie après fixation des canaux, des séries et de la normalisation ;
 \item \(\mathfrak S=\langle90,120\rangle\), semi-groupe infini de hauteurs
       spectrales ;
 \item \(\varepsilon_{\rm res}^{\circ}\), résidu d'une identité angulaire
       exacte, non une extraction finie de l'un quelconque des cinq objets
       précédents.
\end{enumerate}

La coïncidence partielle des indices n'établit pas l'égalité des opérateurs,
des domaines, des codomaines ni de la force probante. En particulier,
\(K_{6\to8}\) utilise deux sous-queues de la hiérarchie, mais n'est pas la hiérarchie
complète ; et aucune de ces sous-queues ne forme un bloc indépendant de la loi
des masses.

\begin{remark}[Séparation canonique des objets temporels]
La notation distingue le registre orienté \(\mathcal R_{12}\), le
calendrier observable \(\Theta_{108}\), la rotation sectorielle \(C_{12}\) et
l'orbite de transport \(\Gamma_{108}\), chacun avec son domaine, son
codomaine et sa loi de composition.
\end{remark}
